% ============================================================================
% CAPITOLO TEMPORANEO — REGISTRO DEI PUNTI APERTI
% ----------------------------------------------------------------------------
% Raccoglie in un unico luogo TUTTE le decisioni progettuali ancora aperte e i
% contenuti da completare disseminati nel documento, così da poterli discutere
% e chiudere in blocco. È TEMPORANEO: a chiusura di ciascuna voce la soluzione
% va riportata nella sezione d'origine (indicata dal rimando) e la voce va
% rimossa da qui. Quando il registro è vuoto, eliminare questo file e il
% relativo \input in DOC_Algisi.tex.
% Ogni voce ha un codice (PA-n / TB-n) per riferirsi ad essa in discussione.
% ============================================================================

\chapter{Punti aperti (\emph{work in progress})}
\label{cap:punti-aperti}

\textit{Nota.} Capitolo temporaneo di servizio. Riunisce, per comodità di discussione, le decisioni ancora da prendere sparse nel documento; ciascuna rimanda alla sezione in cui vive e dove andrà riportata la soluzione una volta chiusa. Va rimosso a lavoro concluso.

\section{Decisioni progettuali aperte}
\label{sec:pa-decisioni}

\subsection{Requisiti da confermare}
\label{sub:pa-requisiti}

\paragraph{PA-1 — Aggregazione delle PII frammentate.} Va deciso se e in che misura il sistema debba ricondurre frammenti dispersi (in uno o più documenti) a un'unica entità coerente, e con quali criteri verificabili (co-occorrenza, similarità semantica, identificatori condivisi). Il requisito non è al momento incluso nell'analisi dei requisiti (cfr.~Capitolo~\ref{cap:requisiti}): la decisione resta aperta.

\paragraph{PA-2 — Costruzione di una knowledge base delle entità.} Va deciso se il sistema debba costruire, a partire dalle PII estratte, una rappresentazione strutturata delle entità (via NER) come base per il confronto col ROPA, e con quale livello di dettaglio. Il requisito non è al momento incluso nell'analisi dei requisiti (cfr.~Capitolo~\ref{cap:requisiti}): la decisione resta aperta. Impatta direttamente su PA-1 e sullo \emph{scope} complessivo (cfr. rischio di \emph{scope creep}, §\ref{sub:rischio-scope-creep}).

\paragraph{PA-12 — Affidabilità della data del documento.} La verifica della conservazione (cfr.~§\ref{sub:b7-retention}) poggia su una data stimata: i metadati interni del file quando esistono, altrimenti la data di ultima modifica sul file system. Entrambe le fonti hanno limiti noti e non eliminabili con l'approccio attuale --- il \texttt{mtime} viene azzerato da copie massive, migrazioni e ripristini da backup; i metadati interni mancano del tutto in alcuni formati (testo semplice, Word legacy, ODF) e possono essere riscritti da una conversione. Resta da decidere se valga la pena ricavare la data dal corpo del documento (formule del tipo «Milano, 12 marzo 2019»), il che richiederebbe un'euristica da validare con lo stesso rigore adottato per il rilevamento delle PII, oppure se accettare la stima attuale dichiarandone i limiti nel verdetto, come avviene ora.

\subsection{Modello dati del ROPA}
\label{sub:pa-ropa}

Voci già isolate come «punti aperti di progettazione» in §\ref{sub:ropa-punti-aperti}, qui riunite.

\paragraph{PA-3 — Granularità della corrispondenza categorie\,--\,PII.} Va fissato se una categoria di dati del ROPA mappi su un singolo \texttt{pii\_type} o su un insieme. Una voce sbrigativa («dati anagrafici») corrisponde di norma a più categorie canoniche, il che suggerisce una relazione 1\,--\,N da formalizzare.

\paragraph{PA-4 — Allineamento nLPD\,--\,GDPR.} Il modello parte dai campi CNIL (impianto GDPR); vanno annotate le differenze rilevanti dell'art.~12 LPD, in particolare soglie ed eccezioni per le PMI, che incidono su quali trattamenti debbano comparire nel registro.

\paragraph{PA-5 — ROPA non strutturato.} Il modello assume un registro già compilato in forma tabellare. Va deciso se (e quanto) supportare un ROPA fornito come testo libero, restando entro il vincolo di non delegare all'AI la strutturazione complessiva del documento.

\subsection{Persistenza del registro PII}
\label{sub:pa-persistenza}

\emph{Risolto per rinvio.} Il tracciamento temporale delle PII (changelog per singola PII) è stato progettato ma \textbf{non realizzato} in questa iterazione: il registro conserva lo stato corrente (§\ref{sec:persistenza-registro}) e l'analisi degli approcci è confluita fra gli sviluppi futuri (§\ref{sec:sf-tracciamento-temporale}). Le decisioni aperte qui sotto (snapshot/event-logger, PA-6/7/8) sono perciò chiuse: si vedano gli approcci A/B/C e il modello dati in §\ref{sec:sf-tracciamento-temporale}.

% ---- RISOLTO PER RINVIO (2026-08-20): feature temporale spostata in §10, decisioni chiuse.
% Testo conservato come traccia storica; da eliminare a pulizia del capitolo.
% \paragraph{Snapshot o event-logger} Da decidere se il metodo per salvare le differenze tra versioni del FS analizzato seguano il paradigma degli snapshot o del'event-logger.
%
% \paragraph{PA-6 — Registrazione della «scomparsa» di una PII: esplicita o implicita?} Quando una scansione non ritrova un'istanza confermata in precedenza, va deciso se scrivere attivamente una \texttt{PIIChange} di tipo \texttt{REMOVED} (esplicita, log autoesplicativo ma con confronto attivo a ogni scansione) o dedurre lo stato «rimossa» a posteriori (implicita, più semplice ma non autosufficiente). Ipotesi preferita: esplicita.
%
% \paragraph{PA-7 — Significato dell'evento «modifica».} Poiché il registro non salva i valori delle PII, un evento di «modifica» può riferirsi solo al cambio di \textbf{posizione} (\texttt{MOVED}) o a una \textbf{riclassificazione} del tipo. Va confermata l'ipotesi di non prevedere un generico \texttt{MODIFIED}, ma di scomporlo in variazioni dal significato preciso.
%
% \paragraph{PA-8 — Vocabolario chiuso dei tipi di variazione.} Si propone di formalizzare \texttt{change\_type} come enumerazione chiusa (\texttt{NEW}, \texttt{CONFIRMED}, \texttt{MOVED}, \texttt{REMOVED}). Resta da confermare se il set è completo per tutti i casi d'uso previsti (dipende anche da PA-6).
% ---- fine blocco risolto ----

\subsection{Scelte implementative}
\label{sub:pa-implementazione}

\paragraph{PA-9 — Tecnologia di persistenza.} Il modello concettuale è \emph{technology-agnostic}; va scelta la tecnologia concreta tra base dati relazionale (es. SQLite, con \emph{join} nativi) e persistenza su file semi-strutturati (JSON/JSONL, più leggera ma senza garanzie relazionali) (cfr.~§\ref{sec:implementazione-persistenza-tecnologia}).

\paragraph{PA-10 — Riuso di Microsoft Presidio come livello di detection.} Va deciso, su base empirica, se costruire il livello regex/NER sopra Presidio (incapsulato dietro il contratto \texttt{PIIDetector}) anziché scriverlo da zero. Il giudizio è affidato a una misurazione diretta (precision/recall/F1) sullo stesso corpus della baseline (cfr.~§\ref{sec:valutazione-presidio}).

\subsection{Metodologia di detection}
\label{sub:pa-metodologia}

\paragraph{PA-11 — Strategia per il modello NER.} Va scelto e motivato l'approccio tra: (A)~zero-shot puro (GLiNER base), (B)~modello pre-fine-tunato da terzi (es. \texttt{gliner-pii}), (C)~fine-tuning leggero in proprio. Dai vincoli di progetto (\emph{zero-training}, assenza di dati annotati, esecuzione locale) l'opzione~C è da scartare a favore di A/B; manca la discussione strutturata con tabella di confronto nel Capitolo~\ref{cap:architettura}.

\section{Contenuti ancora da redigere}
\label{sec:pa-contenuti}

Non decisioni, ma segnaposto e sezioni incomplete sparse nel testo, raccolti qui per non perderli.

\begin{itemize}
	\item \textbf{TB-1} — Architettura: confronto strutturato e scelta del modello NER (tabella dei tre approcci di PA-11) da inserire nel Capitolo~\ref{cap:architettura}.
	% TB-2 RISOLTO (2026-08-20): l'ER è prodotto — stato corrente in Fig.~\ref{fig:er-registro-pii} (§\ref{sub:persistenza-modello-dati}); il modello esteso con \texttt{PIIChange} (design rinviato) in Fig.~\ref{fig:er-changelog-futuro} (§\ref{sec:sf-tracciamento-temporale}). Voce chiusa.
	\item \textbf{TB-3} — Implementazione: roadmap implementativa e di validazione (baseline zero-shot GLiNER, integrazione modello pre-fine-tunato, eventuale esplorazione di fine-tuning) nel Capitolo~\ref{cap:implementazione}.
\end{itemize}
